\documentclass[a4paper,dvips]{slides}
\usepackage[francais,english]{babel}
\usepackage{ifthen}
\usepackage{epsfig}
\usepackage{xspace}

\usepackage{Sty/fig-cap}
\usepackage{Sty/simple}

%%\onlyslides{1-999}

\renewcommand{\arule}[1]{\emph{\mbox{#1}}} % CA rule
\renewcommand{\asquare}[1]{\ensuremath{#1^{2}}}
\newcommand{\agen}[1]{gen.\,#1}
\newcommand{\agencnt}[1]{#1\,gen.}
\newcommand{\aplane}[1]{bit #1}

\newcommand{\mytitle}[1]{\begin{center}\begin{large}#1\end{large}\end{center}}
\newcommand{\mystitle}[1]{\begin{center}#1\end{center}}

\sloppy
\selectlanguage{francais}

\begin{document}

\addtolength{\textheight}{\topmargin}
\addtolength{\textheight}{\headheight}
\addtolength{\textheight}{\headsep}
\addtolength{\textheight}{\footskip}

\setlength{\topmargin}{0pt}
\setlength{\headheight}{0pt}
\setlength{\headsep}{0pt}
\setlength{\footskip}{0pt}

\begin{center}
  \begin{Large}
    Exploration et visualisation\\
    d'espaces du calcul cellulaire
  \end{Large}
  \bigskip
  \begin{large}
    Thierry Delamare\\
  \end{large}
  \bigskip
  \bigskip
  DEA d'intelligence artificielle\\
  Universit\'e Paris-8 et Paris-13\\
  \bigskip
  \today
  \bigskip
\end{center}
\bigskip
\bigskip
\bigskip
\mySfigure{twothirds}{dump}{qmean-9-384}%
{\arule{Qmean} \asquare{512} \agen{384} \aplane{3}}

\begin{slide}
  \begin{center}
    \begin{Large}
      Th\`emes de l'expos\'e\,:
    \end{Large}
  \end{center}
  
  R\`egles et espaces de r\`egles.
  
  Moteurs de calcul et architectures SIMD.
  
  Explorations et visualisations de trajectoires.
\end{slide}

\begin{slide}
  \mytitle{R\`egles\,:}
  Additives, \`a voisinage de Moore~(2D).
  \begin{tabbing}
    XXXXXX\=\kill
    \arule{Life$_1$}\>\texttt{00U100000}\\
    \arule{Anneal$_1$}\>\texttt{0000101111\scriptsize /* annealed majority */}\\
    \arule{Qmean$_8$}\>\texttt{mask = 3; match = 0; delta = -1;}\\
  \end{tabbing}
  \begin{scriptsize}
    \vspace{-\bigskipamount}
    \begin{tabbing}
      XX\=\kill
      \>\texttt{/* Qmean secondary table computation */}\\
      \>\texttt{/* Construct Mean */}\\
      \>\texttt{m = (SUM + C) / 9;}\\
      \>\texttt{/* Destruct it */}\\
      \>\texttt{return~m~\&~mask~==~match~?~m~:~m~+~delta;}
    \end{tabbing}
  \end{scriptsize}
  Congifurations initiales al\'eatoires.
  \mySthreefigure{full}{dump}{hor}{life-00}{anneal-00}{qmean-9-300}%
  {\begin{scriptsize}
    \arule{Life} \asquare{128} \agen{128},
    \arule{Anneal} \asquare{256} \agen{256},\\
    \arule{Qmean} \asquare{512} \agen{300} \aplane{3}
    \end{scriptsize}}
\end{slide}

\begin{slide}
  \mytitle{Moteurs\,:}
  \mystitle{Lookup-tables, pipelines, voisinages virtuels}
  \myStwofigureA{full}{fig}{hor}{apply-1}{func-4}%
  {Pipeline d'application de \arule{Life} et \arule{Anneal}\\
    et deux tables non disjointes}
  \mySfigure{full}{fig}{func-7}%
  {Fonction d'application du voisinage de \arule{Qmean}}
\end{slide}

\begin{slide}
  \mytitle{Explorations\,:}
  \mystitle{Invariants, limites}
  \mySfigure{full}{plot}{anneal-04}%
  {\arule{Anneal} $2^{{7\cdots10}^2}$ $3500\cdots110000$ gen.}
  \mySfigureR{full}{.9}{plot}{anneal-02}%
  {\arule{Anneal} \asquare{1024}}
\end{slide}

\begin{slide}
  \mytitle{Visualisations\,:}
  \mystitle{Extractions, diff\'erences, accumulations}
  \mySfigure{full}{dump}{life-time-xor}%
  {\arule{Life} \asquare{64} \agencnt{6}}
  \myStwofigure{full}{plot}{hor}{lifex-sum-1}{lifex-sum-2}%
  {\arule{Life} \asquare{64} \agencnt{512}}
\end{slide}

\begin{slide}
  \mystitle{Espace-temps}
  \mySfigure{full}{dump}{life-3d-New}%
  {\arule{Life} \asquare{64} \agencnt{512}}
\end{slide}

\begin{slide}
  \mystitle{Projections relativistes, tranches, transparences, iso-surfaces}
  \mySfigure{full}{dump}{slide-2}{%
    \arule{Anneal} \acube{256}, %
    \arule{Life} \acube{256},\\
    \arule{Qmean} \acube{128}, %
    \arule{Anneal} \asquare{64} \agencnt{400}}
\end{slide}

\begin{slide}
  \mytitle{Espaces de r\`egles\,:}
  \begin{footnotesize}
    Soit une r\`egle $R$ d'un AC de $k = 2^n$ \'etats,\\
    $d$ dimensions et de rayon de voisinage $r$.\\
    $N = k^{dr+1} = 2^b$ est le nombre d'entr\'ees de la table.
  \end{footnotesize}
  \begin{eqnarray}
    \lambda & = & (N-\#0)/N \nonumber \\
    {\rm flip} & = & \frac{1}{N}\sum_{i=0}^{N-1}{\rm
      diff}(R_{i}-R_{(i+1){\rm mod}N}) \nonumber \\
    {\rm diff}(x) & = & \left\lbrace
    \begin{array}{l}
      1\ {\rm if}\ x \ne 0\\ 0\ {\rm if}\ x = 0
    \end{array} \right. \nonumber
  \end{eqnarray}
  \begin{footnotesize}
    Pour un AC de $k = 2$ \'etats les $b = dr+1$ bits du voisinage donnent
    un vecteur $E$ de $b$ valeurs.
  \end{footnotesize}
  \begin{eqnarray}
    E_{b}& = &\frac{1}{N} \sum_{i = 0}^{N/2 - 1} R_{e_{ib}} -
    R_{e_{ib} + 2^b} \nonumber \\
    e_{ib}& = &2i + i \bmod 2^b \nonumber \\
    \hat{E}& = &\frac{1}{b} \sum_{i = 0}^{b} E_b \nonumber
  \end{eqnarray}
\end{slide}

\begin{slide}
  \mystitle{Valeurs de \alambda et de l'ench\^assement pour les
    $\frac{0\cdots 500}{59049}$ AC additifs \texttt{01U}}
  \mySfigure{full}{plot}{E-lambda-1}%
  {Valeurs de \protect\alambda et $E$}
  \begin{eqnarray}
    E_{life} & = & {0.617\cdots 0.890\cdots 0.617} \nonumber \\
    \hat{E}_{life} & = & 0.647569 \nonumber \\
    \lambda_{life} & = & 0.273438 \nonumber
  \end{eqnarray}
\end{slide}

\begin{slide}
  \mytitle{SIMD\,:}
  Transformation d'architecture SIMD en table de transition \`a
  voisinage virtuel pour la m\'e\-moire et les instructions.
  \myStwofigureA{full}{fig}{hor}{gap-pe}{gap}%
  {Structure d'un PE GAPP\\
    et fonction d'application de la table \arule{Gapp}}
  Transformation de table de transition en DFA pour
  interpr\'etation par algorithme SIMD.
  \myStwofigureA{full}{fig}{hor}{life-tree}{life-dfa}%
  {Une partie de l'arbre de d\'ecodage de \arule{Life}\\
    et Le DFA de \arule{Life}}
\end{slide}

\begin{slide}
  \begin{center}
    \begin{Large}
      Conclusion
    \end{Large}
  \end{center}
  \mytitle{Propositions pour une \'etude pratique des espaces du calcul
    cellulaire\,:}

  Un usage g\'en\'eralis\'e de table de transitions.

  Un couplage des architectures SIMD et des tables de transitions.

  Une mesure de l'ench\^assement des tables de transitions.

  Des m\'ethodes de visualisations d'espace-temps.
\end{slide}

%%
%% TUBE-WORM MOSAIC
%%
\begin{slide}
  \mystitle{Espace de r\`egle:~huit variations sur les param\`etres
    d'un algorithme de production de tables de transition}
  \mySfigure{full}{dump}{tube-worm-mix}%
  {\arule{Tube-Worm} \asquare{256} \agencnt{64}}
\end{slide}

%%
%%
%% ANNEAL 3D
%%
\begin{slide}
  \mystitle{Fronti\`eres et textures}
  \mySfigure{full}{dump}{anneal-3d}{\arule{Anneal} \acube{512}}
\end{slide}

%%% XU
%%
\begin{slide}
  \mystitle{Un AC algorithmique}
  \mySfigure{full}{dump}{xu-photo-1}%
  {\arule{Xu} \acube{128}}%
\end{slide}

\end{document}

%% Local Variables: 
%% mode: auto-fill
%% iso-accents-minor-mode: nil
%% TeX-master: t
%% Local IspellParsing: "latex-mode ~latin1"
%% Local IspellPersDict: "~/.ispell_lisp"
%% LocalWords: "paper dvips slides francais english babel ifthen epsfig xspace"
%% LocalWords: "Sty fig-cap gen mynote pt Thierry Delamare DEA twothirds dump XX"
%% LocalWords: "qmean SIMD XXXXXX Anneal mask secondary computation Construct it"
%% LocalWords: "Mean SUM Destruct return hor anneal apply func santes lifex-sum"
%% LocalWords: "life-time-xor lifex-sum Espace-temps d-New relativiste rel tsum"
%% LocalWords: "relmov d'espace-temps anneal-sum iso-surface tsum-ray-one tsum"
%% LocalWords: "trans-byte trans-bits diag E-lambda tube-worm-mix Tube-Worm PE"
%% LocalWords: "gap-pe GAPP gap gapp life-tree life-dfa DFA cellulaire-SIMD l'AC"
%% LocalWords: "xu-photo xu xu-Cfilm \'equi"
%% End: 
